%!TeX root = ../principal.tex
%!TeX encoding = utf-8
\chapter{Método}
\label{metodo}
\section{Busca}
Como base de dados foram usados IEEE Xplore, ACM Digital Library, Scopus, SpringerLink, and Wiley Online Library.
Em cada uma dessas bases foi executada a seguinte string de busca:
\begin{verbatim}
(
    ("forest smoke" OR "forest fire")
    AND ("detection" OR "identification")
    AND ("image*" OR "video*" OR "vision*")
    AND NOT "survey" AND NOT "review"
    AND ("AI" OR "ML" OR "machine learning")
)
\end{verbatim}
Foram escolhidos apenas artigos de 2020-2025.
Foram excluídos artigos que não podem ser acessados gratuitamente.
Artigos encontrados:
\begin{table}[h]
\centering
\begin{tabular}{ |l c| }
    \hline
    IEEE Xplore &  8 \\
    \hline
    ACM &  9 \\
    \hline
    Scopus & 23 \\
    \hline
    SpringerLink & 51 \\
    \hline
    Wiley & 20 \\
    \hline
\end{tabular}
\caption{Artigos encontrados pela string de busca}
\label{found_papers}
\end{table}

Em seguida foram classificados. Em tipo de pesquisa (secundária ou primária), se é sobre Detecção de incêndio florestal usando inteligência artificial (sim, não, null), tipo de imagem (estacionária, drone, satélite, video, null) e.
\begin{table}[h]
\centering
\begin{tabular}{ |l c| }
    \hline
    Primária & 101 \\
    \hline
    Secundária & 10 \\
    \hline
\end{tabular}
\caption{Tipo de pesquisa}
\label{research_type}
\end{table}

\begin{table}[h]
\centering
\begin{tabular}{ |l c| }
    \hline
    Sim & 28 \\
    \hline
    Não & 74 \\
    \hline
    NULL & 9 \\
    \hline
\end{tabular}
\caption{É sobre Detecção de incêndio florestal usando inteligência artificial, nulos foram eliminados por serem pesquisa secundária antes de serem classificados nesta categoria.}
\label{isaboutforestfire}
\end{table}

\begin{table}[h]
\centering
\begin{tabular}{ |l c| }
    \hline
    drone & 9 \\
    \hline
    estacionária & 10 \\
    \hline
    satélite & 8 \\
    \hline
    video & 10 \\
    \hline
    NULL & 74 \\
    \hline
\end{tabular}
\caption{Tipo de imagem, nulos foram eliminados pelas categorias anteriores antes de serem classificados nesta.}
\label{isaboutforestfire}
\end{table}

Foram excluídos os artigos que são pesquisa secundária, não são sobre detecção de incêndio florestal usando inteligência artificial e usam tipo de imagem diferente de drone ou estacionária. Após isso sobraram 16 artigos, desses foram escolhidos um de cada de tipo de IA dando prioridade aqueles que providenciam códigos e dados. Finalmente sobrando 10 artigos.

\section{Treinamento}
Treinamos cada modelo por aproximadamente 10 horas usando o Dataset D-Fire \citar{dfire}, Foi usado o loss inicial fornecido em cada paper e um decaimento de loss usando uma curva sigmoide seguindo a quantidade de épocas nos respectivos papers.

Do dataset usamos a divisão de traino (X imagens), validação (X imagens) e teste (X imagens) fornecidos pelo próprio dataset. Para o TFNet foram excluídas todas as imagens (X imagens) que não contém objetos, pois imagens sem objetos atrapalham imensamente o treinamento de models baseados no RT-DETR.
